\documentclass[11pt,a4paper]{article}
\usepackage[top=2.8cm, bottom=2.8cm, left=3cm, right=3cm]{geometry}
\usepackage[T1]{fontenc}
\usepackage[utf8]{inputenc}
\usepackage{lmodern}
\usepackage{microtype}
\usepackage[colorlinks=true, linkcolor=linkblue, urlcolor=accent,
            pdftitle={L1-04 Load Balancing -- System Design Foundations},
            bookmarks=true]{hyperref}
\usepackage{xcolor}
\usepackage{titlesec}
\usepackage{enumitem}
\usepackage{booktabs}
\usepackage{tabularx}
\usepackage{tcolorbox}
\usepackage{fancyhdr}
\usepackage{amsmath}
\usepackage{parskip}
\usepackage{mdframed}
\usepackage{graphicx}
\usepackage{tikz}
\usepackage{setspace}
\usepackage{soul}
\usepackage{array}
\usepackage{longtable}

\definecolor{primary}{HTML}{1A1A2E}
\definecolor{accent}{HTML}{C0392B}
\definecolor{highlight}{HTML}{1A5276}
\definecolor{linkblue}{HTML}{154360}
\definecolor{lightbg}{HTML}{F4F6F7}
\definecolor{quizbg}{HTML}{EBF5FB}
\definecolor{termbg}{HTML}{EAF2FF}
\definecolor{curiobg}{HTML}{F5EEF8}
\definecolor{greenbg}{HTML}{EAFAF1}
\definecolor{notebg}{HTML}{FDF2E9}
\definecolor{accentlight}{HTML}{FADBD8}

\titleformat{\section}{\large\bfseries\color{highlight}}{\thesection.}{0.6em}{}[\vspace{2pt}\color{highlight!60}\rule{\linewidth}{0.6pt}]
\titleformat{\subsection}{\normalsize\bfseries\color{primary}}{\thesubsection}{0.5em}{}
\titleformat{\subsubsection}{\normalsize\itshape\color{primary!80}}{}{0em}{}
\titlespacing{\section}{0pt}{18pt}{8pt}
\titlespacing{\subsection}{0pt}{12pt}{4pt}

\pagestyle{fancy}
\fancyhf{}
\fancyhead[L]{\small\color{highlight}\textbf{L1-04 --- Load Balancing}}
\fancyhead[R]{\small\color{primary!70}System Design Interview Prep}
\fancyfoot[C]{\small\color{primary!60}\thepage}
\renewcommand{\headrulewidth}{0.4pt}
\renewcommand{\headrule}{\hbox to\headwidth{\color{highlight!40}\leaders\hrule height \headrulewidth\hfill}}

\tcbuselibrary{skins, breakable, hooks}
\newtcolorbox{termbox}[1]{enhanced,breakable,colback=termbg,colframe=highlight,coltitle=white,fonttitle=\bfseries\small,title={#1},left=8pt,right=8pt,top=6pt,bottom=6pt,attach boxed title to top left={yshift=-2.5mm,xshift=6mm},boxed title style={colback=highlight,sharp corners,left=4pt,right=4pt}}
\newtcolorbox{industrybox}[1]{enhanced,breakable,colback=greenbg,colframe=highlight!50!black,coltitle=white,fonttitle=\bfseries\small,title={Industry Data: #1},left=8pt,right=8pt,top=6pt,bottom=6pt,attach boxed title to top left={yshift=-2.5mm,xshift=6mm},boxed title style={colback=highlight!60!black,sharp corners,left=4pt,right=4pt}}
\newtcolorbox{trapbox}[1]{enhanced,colback=accentlight,colframe=accent,coltitle=white,fonttitle=\bfseries\small,title={Interview Trap: #1},left=8pt,right=8pt,top=6pt,bottom=6pt,attach boxed title to top left={yshift=-2.5mm,xshift=6mm},boxed title style={colback=accent,sharp corners,left=4pt,right=4pt}}
\newtcolorbox{thinkbox}{enhanced,colback=notebg,colframe=accent!60,left=8pt,right=8pt,top=6pt,bottom=6pt,leftrule=3pt,rightrule=0pt,toprule=0pt,bottomrule=0pt}
\newtcolorbox{curiobox}[1]{enhanced,breakable,colback=curiobg,colframe=primary!60,coltitle=white,fonttitle=\bfseries\small,title={Q: #1},left=8pt,right=8pt,top=6pt,bottom=6pt,attach boxed title to top left={yshift=-2.5mm,xshift=6mm},boxed title style={colback=primary!70,sharp corners,left=4pt,right=4pt}}
\newtcolorbox{quizbox}{enhanced,breakable,colback=quizbg,colframe=highlight,coltitle=white,fonttitle=\bfseries,title={Self-Quiz},left=10pt,right=10pt,top=8pt,bottom=8pt,attach boxed title to top left={yshift=-2.5mm,xshift=6mm},boxed title style={colback=highlight,sharp corners,left=4pt,right=4pt}}
\newtcolorbox{answerbox}{enhanced,breakable,colback=lightbg,colframe=primary!40,coltitle=white,fonttitle=\bfseries\small,title={Model Answers},left=10pt,right=10pt,top=8pt,bottom=8pt,attach boxed title to top left={yshift=-2.5mm,xshift=6mm},boxed title style={colback=primary!60,sharp corners,left=4pt,right=4pt}}
\newtcolorbox{insightbox}{enhanced,colback=primary!5,colframe=primary,left=8pt,right=8pt,top=6pt,bottom=6pt,leftrule=4pt,rightrule=0.4pt,toprule=0.4pt,bottomrule=0.4pt}

\setstretch{1.15}

% ─────────────────────────────────────────────────────────────────────────────
\begin{document}

% ── Title Page ───────────────────────────────────────────────────────────────
\begin{titlepage}
  \thispagestyle{empty}
  \vspace*{2cm}
  \begin{center}
    {\small\color{highlight}\textbf{LAYER 1 --- FOUNDATIONS}}\par
    \vspace{0.4cm}
    {\small\color{primary!60}Lesson L1-04}\par
    \vspace{1.2cm}
    {\LARGE\bfseries\color{primary} Load Balancing}\par
    \vspace{0.5cm}
    {\large\color{primary!70} Distributing Traffic Across a Fleet of Servers}\par
    \vspace{0.4cm}
    {\small\color{primary!50} System Design Interview Prep --- 2025 Edition}\par
    \vspace{2.2cm}

    \begin{insightbox}
      \textbf{Why This Lesson Matters.} Every large-scale system you will be asked to design in
      an interview --- a notification pipeline, a URL shortener, a social feed, a payment processor
      --- implicitly assumes that you know how to route traffic across many servers without making
      any single server the bottleneck or single point of failure. Load balancing is the invisible
      skeleton of every distributed system. Interviewers frequently escalate from ``design a login
      service'' to ``now scale it to 100 million users'' within two minutes. Your answer requires
      a load balancer at its foundation. This lesson makes that foundation explicit.
    \end{insightbox}

    \vspace{1.4cm}
    {\small\color{primary!50}
      \textbf{Data sources:} AWS Elastic Load Balancing documentation (2024--2025);
      HAProxy benchmarks on AWS Graviton (HAProxy, Inc., 2024);
      Cloudflare network infrastructure documentation (2024);
      Apache Cassandra architecture documentation (2024);
      Netflix OSS: Eureka and Ribbon (2020--2024);
      Google Cloud Load Balancing documentation (2024--2025);
      VRRP / Keepalived HA documentation (2024).
    }\par
  \end{center}
\end{titlepage}

% ── Table of Contents ────────────────────────────────────────────────────────
\newpage
\tableofcontents
\newpage

% ─────────────────────────────────────────────────────────────────────────────
\section{Why This Exists: The Problem of One Server}

Early web services were simple: one machine, one IP address, users connect directly. That model
collapses the moment you try to scale. A single physical server has a fixed ceiling on the number
of simultaneous TCP connections it can hold open, a fixed amount of CPU to process requests, and a
finite network interface to push data through. More critically, that single machine is one power
supply failure, one kernel panic, or one deployment gone wrong away from total outage. The question
is not \emph{if} a server will fail --- it is \emph{when}, and whether your users notice.

The instinctive first response is to buy a bigger server. This is called \emph{vertical scaling} or
``scaling up.'' It works, up to a point. But the largest commercially available machines top out at
a few thousand cores and a few hundred terabytes of RAM, and they cost tens of millions of dollars.
More practically, even the most powerful single server is still a single point of failure. No
vertical scaling strategy survives the data-center power outage that hit AWS us-east-1 in December
2021 and brought down thousands of services simultaneously.

The enduring answer is \emph{horizontal scaling}: instead of one large server, operate many
commodity servers in parallel. That simple idea unlocks essentially unlimited capacity --- you can
keep adding machines as traffic grows. But horizontal scaling immediately raises a new question:
\emph{how do clients find the right server?} If you have fifty backend instances, each with its own
IP address, clients cannot be expected to pick one at random and hope for the best. You need
something that sits in front of the fleet, accepts every incoming request, and makes an intelligent
decision about which backend should handle it. That something is a \textbf{load balancer}.

A load balancer solves three distinct problems at once. First, it \emph{distributes traffic}
so that no single backend becomes a hot spot while others sit idle. Second, it provides
\emph{health-based routing}: when a backend fails its health check, the load balancer stops sending
it traffic and reroutes requests to the surviving instances, making failures invisible to the client.
Third, it presents a \emph{single, stable endpoint} to the outside world --- one IP address or
hostname --- regardless of how many backends exist behind it or how frequently they change.

\begin{insightbox}
  \textbf{Mental model.} Think of a load balancer as an air-traffic control tower. Planes
  (requests) arrive continuously. The tower (LB) knows which runways (backends) are open,
  which are congested, and which are temporarily closed for maintenance. It directs each
  plane to the best available runway. Planes never talk to runways directly; they always go
  through the tower.
\end{insightbox}

% ─────────────────────────────────────────────────────────────────────────────
\section{Layer 4 vs.\ Layer 7: What the Load Balancer Can See}

Before deciding \emph{how} to distribute traffic, a load balancer must decide \emph{where} in the
network stack it operates. This is the most conceptually important distinction in the topic, and the
one interviewers probe most frequently.

\subsection{Layer 4 (Transport Layer) Load Balancing}

A Layer 4 (L4) load balancer operates at the TCP/UDP transport layer. It makes routing decisions
based entirely on the network tuple: source IP, destination IP, source port, destination port, and
protocol. It does not read any of the payload. When a TCP connection arrives, the L4 balancer picks
a backend and then forwards all subsequent packets in that connection to the same backend. From the
backend's perspective, the connection looks like it came from the client (or, in many implementations,
from the load balancer's own IP via SNAT --- source network address translation).

Because an L4 balancer never decrypts or parses the payload, it is extremely fast. It can operate
at line rate on commodity hardware. It also works for any protocol built on top of TCP or UDP:
HTTP, gRPC, MySQL, Redis, SMTP, custom binary protocols --- the L4 balancer does not care. This
generality is both its strength and its limitation.

\begin{termbox}{Layer 4 Load Balancer}
  Routes connections based on the network 4-tuple (src IP, src port, dst IP, dst port). Sees no
  application-layer data. Operates at near wire-speed. Cannot inspect HTTP headers, cookies, or
  request paths. Suitable for any TCP/UDP protocol.
\end{termbox}

\subsection{Layer 7 (Application Layer) Load Balancing}

A Layer 7 (L7) load balancer terminates the connection from the client --- including TLS decryption
if the connection is HTTPS --- and then reads the actual HTTP request before forwarding it to a
backend. Because it can see the full application message, it can make routing decisions based on:
the URL path (\texttt{/api/v2/} goes to one pool, \texttt{/static/} goes to another), HTTP headers
(route \texttt{X-Region: eu} to the EU cluster), cookies (send users with a specific session cookie
to the same backend), or even the body of the request.

This richness comes at a cost. Terminating TLS and parsing HTTP requires significantly more CPU than
simple packet forwarding. An L7 balancer introduces a small but non-zero latency overhead. And
because it terminates the client connection, it establishes a new, separate connection to the
backend, meaning backends see the load balancer's IP rather than the original client's IP unless the
balancer injects an \texttt{X-Forwarded-For} header.

\begin{termbox}{Layer 7 Load Balancer}
  Terminates the client TCP connection, reads and parses the HTTP/HTTPS request, then opens a new
  connection to a chosen backend. Can route on URL path, headers, cookies, or query parameters.
  Handles TLS termination. Adds modest CPU overhead. Backends see the LB's IP unless
  \texttt{X-Forwarded-For} is configured.
\end{termbox}

\begin{thinkbox}
  \textbf{Why terminate TLS at the load balancer rather than at each backend?} The private key
  required for TLS decryption is highly sensitive. Distributing it across dozens of backend servers
  multiplies your exposure surface. Centralising TLS termination at the load balancer means only
  one (or a small HA pair) of machines holds that key, simplifying certificate management and
  rotation. Backends then communicate over the internal network using plain HTTP or mTLS with an
  internal CA --- a much simpler key management story.
\end{thinkbox}

\subsection{Choosing Between L4 and L7}

The practical decision rule is straightforward. If you need content-based routing, cookie-based
session affinity, A/B traffic splitting, WebSocket upgrade handling, or gRPC-specific load
balancing, you need L7. If you need the absolute lowest latency and highest throughput for a
non-HTTP protocol --- a real-time multiplayer game server, a database proxy, a raw socket streaming
service --- L4 is the right choice. Many production stacks layer both: an L4 balancer at the
internet edge for raw TLS offload and DDoS absorption, feeding into L7 balancers for per-service
content routing.

% ─────────────────────────────────────────────────────────────────────────────
\section{Balancing Algorithms: How to Pick a Backend}

Given a pool of healthy backends, the load balancer must implement an algorithm to decide which one
receives the next request. The choice of algorithm has significant implications for fairness,
latency, and hot-spot avoidance.

\subsection{Round Robin}

Round Robin is the simplest algorithm: the load balancer maintains a list of backends and cycles
through them in order. Request 1 goes to backend A, request 2 to B, request 3 to C, request 4
back to A, and so on. It is stateless, trivially fast, and distributes requests evenly when all
requests take approximately the same time to process and all backends have equal capacity.

Round Robin breaks down when requests are heterogeneous. If some requests take 10ms and others
take 2,000ms (a slow database query, a video transcode), the raw count of requests per backend is
a poor proxy for actual load. A backend that received several expensive requests will be heavily
loaded even though it has received the same number of requests as a lightly loaded peer.

\subsection{Weighted Round Robin}

Weighted Round Robin assigns each backend a numeric weight reflecting its relative capacity. A
backend with weight 4 receives four requests for every one request sent to a backend with weight 1.
This is the right choice when your fleet is heterogeneous --- perhaps you are gradually migrating
from one instance type to another, or you are canary-testing a new version with 5\% of traffic.
AWS ALB supports target group weights natively, which is exactly how engineers at Amazon perform
weighted traffic splits for deployments.

\subsection{Least Connections}

Least Connections routes each new request to the backend currently handling the fewest active
connections. It is inherently adaptive: if a backend is slow and its connection queue grows, it
automatically receives less new traffic until it catches up. This makes it significantly better
than Round Robin for workloads with variable request duration, such as database queries, file
uploads, or any endpoint with long-tail latency.

\subsection{Least Response Time}

Least Response Time is an extension of Least Connections that incorporates actual backend latency:
the load balancer tracks the average response time of each backend (via a rolling window) and routes
new requests to the backend with the lowest combination of active connections and response time.
HAProxy and NGINX Plus both implement variants of this. It is the highest-fidelity algorithm for
adaptive load distribution but requires the load balancer to maintain per-backend latency state,
which adds overhead.

\subsection{IP Hash}

IP Hash computes a hash of the client's source IP address and maps that hash to a specific backend.
The same client IP always lands on the same backend, providing a form of server affinity without
explicit session tracking. It is deterministic and stateless from the load balancer's perspective.
Its weakness: it cannot rebalance. If one backend becomes slow, clients whose IP hashes to that
server are stuck there. And if the backend set changes (a server is added or removed), almost all
existing hash mappings are invalidated, causing a complete redistribution of sessions --- exactly
the problem that consistent hashing was designed to solve.

\subsection{Consistent Hashing}

Consistent hashing is the intellectually richest algorithm in this list, and the one interviewers
probe most deeply because it appears everywhere in distributed systems: not just load balancers, but
databases, CDN edge routing, and distributed caches.

The core insight is this: with naive modular hashing (\texttt{key \% N} where \texttt{N} is the
number of backends), adding or removing even one backend invalidates the bucket assignment for
nearly every key. Consistent hashing solves this by imagining the hash space as a ring from 0 to
$2^{64}-1$. Each backend is hashed to one or more positions on that ring. To route a request, you
hash the routing key (client IP, session ID, or request key) to a position on the ring, then walk
clockwise until you hit a backend node. That backend owns the request.

The elegant property: when you add a new backend, it takes over only the segment of the ring
between itself and its clockwise predecessor. All other request-to-backend assignments remain
unchanged. Similarly, when a backend is removed, only its segment's keys move to the next clockwise
node. In expectation, only $\frac{1}{N}$ of keys are remapped when a server is added or removed,
compared to almost 100\% with modular hashing.

\begin{thinkbox}
  \textbf{But what about hot spots on the ring?} With a small number of physical nodes, the spacing
  between them on the ring is uneven by chance, causing some backends to own disproportionately
  large ring segments. The solution is \emph{virtual nodes} (vnodes): instead of placing each
  physical server once on the ring, you place it many times (at many different hash positions),
  using different seeds. Apache Cassandra defaults to 256 vnodes per physical node. With 256 points
  per server, the ring spacing becomes statistically uniform even with just a few physical servers.
  DynamoDB uses a conceptually identical mechanism internally, though it hides the ring entirely
  behind its partition abstraction.
\end{thinkbox}

\begin{industrybox}{Consistent Hashing in Production --- Cassandra and CDN Edge (2024)}
  Apache Cassandra's architecture document specifies a $2^{64}$ Murmur3 hash ring with 256 virtual
  nodes per physical node by default (as of Cassandra 4.0+). This means a 12-node cluster has
  3,072 positions on the ring, giving statistically even load distribution. Redis Cluster takes a
  different approach: it uses 16,384 fixed hash slots with CRC16, a simpler but less adaptive
  scheme. Cloudflare's Unimog edge load balancer uses consistent hashing to route flows to edge
  servers within a data center, ensuring that adding capacity during a traffic spike does not
  invalidate all existing connection state.
\end{industrybox}

% ─────────────────────────────────────────────────────────────────────────────
\section{Health Checks: How Dead Servers Are Detected}

A load balancer that routes traffic to a dead or degraded backend defeats its own purpose. Health
checking is the mechanism by which a load balancer continuously monitors the state of its backend
pool and removes unhealthy members from rotation.

\subsection{Active Health Checks}

In active health checking, the load balancer periodically probes each backend itself, independent of
real user traffic. The probe can be as simple as a TCP connection attempt (``can I reach port 443?'')
or as sophisticated as an HTTP GET to a dedicated health endpoint (\texttt{/health} or
\texttt{/healthz}) that the backend exposes. That endpoint typically checks whether the application's
own dependencies --- its database connection, its cache, its downstream APIs --- are alive before
returning 200 OK.

The key parameters are the health check interval (how often to probe), the failure threshold (how
many consecutive failures before the backend is marked down), and the recovery threshold (how many
consecutive successes before the backend is returned to rotation). AWS ALB defaults to a 30-second
interval with 5 failures required to mark a target unhealthy --- giving a 2.5-minute window before
a truly dead server is removed. In latency-sensitive production systems, teams often tighten this
to a 5-second interval with 2 failures, trading slightly more health-check traffic for faster
failure detection.

\subsection{Passive Health Checks}

Passive health checking (also called ``outlier detection'' in Envoy's terminology) infers backend
health from the responses to real traffic rather than synthetic probes. If a backend returns 5xx
errors at an elevated rate, or if its response latency spikes beyond a threshold, the load balancer
marks it as degraded and reduces or eliminates its share of traffic. Envoy Proxy, used extensively
by Lyft, Airbnb, and Google internally, implements passive outlier detection as a first-class
feature: after a configurable number of consecutive 5xx responses from a backend, Envoy ejects it
from the load-balancing pool for an exponentially increasing duration.

\begin{industrybox}{AWS ALB Health Check Defaults (AWS, 2024)}
  AWS Application Load Balancer health check defaults: protocol HTTP or HTTPS, path \texttt{/},
  interval 30 seconds, timeout 5 seconds, healthy threshold 5, unhealthy threshold 2. This means
  a backend must fail 2 consecutive checks (60 seconds minimum) before being marked unhealthy, and
  must pass 5 consecutive checks (150 seconds) before being returned to service. Teams running
  high-availability payment infrastructure on AWS typically override these to interval=10s,
  unhealthy threshold=2, healthy threshold=2 --- cutting worst-case failure detection from 60
  seconds to 20 seconds.
\end{industrybox}

% ─────────────────────────────────────────────────────────────────────────────
\section{Session Persistence and the Stateless Design Imperative}

\subsection{What Sticky Sessions Are}

When a user logs into a web application and the server stores session state --- the shopping cart
contents, the authentication token, the preferences --- in the server's local memory (a HashMap, a
Redis client-side cache, or just a plain variable), then that user's \emph{subsequent requests must
reach the same server}. If the load balancer routes a follow-up request to a different backend, the
new backend has no knowledge of the session and the user appears to be logged out.

Session persistence, also called ``sticky sessions'' or ``session affinity,'' is the load balancer
feature that solves this by ensuring a given client always lands on the same backend. The most
common mechanism is a cookie: on the first request, the load balancer sets a cookie (e.g., AWS ALB
uses a cookie named \texttt{AWSALB}) that encodes which backend handled the request. On every
subsequent request, the load balancer reads that cookie and routes accordingly.

\begin{trapbox}{Sticky Sessions as a Design Answer}
  A frequent interview mistake is treating sticky sessions as the correct, scalable solution to
  session state. It is not. Sticky sessions are a band-aid that reintroduces single-server
  coupling at the application layer. When the sticky backend fails, all of that server's sessions
  are lost anyway. When the fleet autoscales, many existing sticky assignments point to the new
  servers and the load balancer must rebuild its mapping tables. The load balancer now carries state
  (the session table), making it harder to replace or restart. The \textbf{correct answer} is to
  design your backends to be stateless: store all session state in a shared, distributed store
  (Redis, DynamoDB, or a database) that any backend can read. The load balancer then has no
  affinity table to maintain and can route freely. When your interviewer asks ``how do you handle
  user sessions?'' the answer is \emph{not} sticky sessions. It is: ``session tokens are stored
  in a distributed cache accessible by all backends.''
\end{trapbox}

\subsection{When Sticky Sessions Are Unavoidable}

Sticky sessions are not always wrong --- just usually wrong. The legitimate use cases are narrow:
long-lived stateful protocols where session state cannot be externalised cheaply (certain WebRTC
signalling flows, some legacy financial trading system protocols) or situations where externalising
state has prohibitive latency cost. Even then, the correct modern approach is to use an L7 load
balancer that supports WebSocket upgrades and proxies the long-lived connection transparently,
avoiding the need for session affinity entirely.

% ─────────────────────────────────────────────────────────────────────────────
\section{Global Server Load Balancing: Routing Across Data Centres}

So far we have described load balancing within a single data centre or region. At global scale ---
when you operate services in multiple geographic regions and want users to reach the nearest healthy
region --- you need a layer above the regional load balancers. This is \textbf{Global Server Load
Balancing (GSLB)}.

\subsection{GeoDNS}

The simplest GSLB mechanism is GeoDNS: when a client resolves your domain name, the authoritative
DNS server inspects the client's IP address, infers its geographic location, and returns the IP
address of the nearest regional load balancer. A user in Tokyo resolving \texttt{api.example.com}
receives the IP of the Tokyo region LB; a user in Frankfurt receives the Frankfurt IP.

GeoDNS is cheap to implement and universally compatible. Its fundamental limitation is that DNS
responses are cached. TTL values of 30--300 seconds mean that if the Tokyo region goes down,
clients that already cached the Tokyo IP will continue attempting to connect to it for up to five
minutes before re-resolving and getting routed to the next region. This makes GeoDNS unsuitable
as a sole failover mechanism when your RTO (recovery time objective) is measured in seconds rather
than minutes.

\subsection{Anycast}

Anycast is a network routing technique where the same IP address is simultaneously announced from
multiple geographic locations. When a client sends a packet to that IP, the internet's BGP routing
infrastructure directs the packet to the \emph{nearest} announcement point --- measured in network
hops, not geographic distance. No DNS manipulation is required; the routing is invisible to the
client and happens at the IP layer.

Anycast delivers sub-millisecond failover at the routing layer. If the Frankfurt point of presence
goes offline, BGP converges and subsequent packets from European clients are rerouted to the next
nearest announcement (Amsterdam or Paris) within seconds, not minutes. This is why every major
CDN and DDoS mitigation provider uses anycast as the foundation of their global network.

\begin{industrybox}{Cloudflare Anycast Network (Cloudflare, 2024)}
  Cloudflare operates its load balancing and CDN across 330+ data centre locations (points of
  presence) interconnected with over 13,000 network peers. Every Cloudflare service --- DNS, load
  balancing, CDN, DDoS scrubbing --- is served from the same anycast IP space. Cloudflare reports
  that its network reaches approximately 95\% of the world's internet-connected population within
  50ms of latency. This means that when Cloudflare's load balancer accepts your HTTPS connection,
  TLS termination and health-check-based routing happen at the edge location closest to the client,
  with the re-encrypted connection forwarded over Cloudflare's private backbone to your origin.
  Google Cloud Load Balancing uses an architecturally similar approach: a single anycast VIP
  programmed globally into Google's Maglev infrastructure, with GFE (Google Front End) edge nodes
  terminating connections at the nearest Google PoP and forwarding internally to the selected
  backend region.
\end{industrybox}

\begin{thinkbox}
  \textbf{Anycast sounds perfect. Why does anyone still use GeoDNS?} Anycast requires you to
  announce your IP block via BGP from each location, which means you must own or lease BGP-capable
  network infrastructure. GeoDNS, by contrast, works with any DNS provider and any IP addressing
  scheme. For teams without their own ASN and BGP infrastructure, GeoDNS plus a short TTL (30--60
  seconds) is the practical choice for geographic routing. Many production systems combine both:
  anycast for the CDN edge layer, GeoDNS for application-tier regional routing.
\end{thinkbox}

% ─────────────────────────────────────────────────────────────────────────────
\section{Load Balancer High Availability: Who Watches the Watcher?}

A load balancer that itself is a single point of failure is self-defeating. Any serious production
deployment requires the load balancer to be highly available.

\subsection{Active--Passive with VRRP}

The most common HA pattern for software load balancers (HAProxy, NGINX) on bare metal or VMs is
active--passive failover via the \textbf{Virtual Router Redundancy Protocol (VRRP)}. Two (or more)
load balancer instances share a \emph{virtual IP address} (VIP). At any moment, one instance --- the
``master'' --- owns the VIP and processes all traffic. The other instances --- the ``backups'' ---
listen for a heartbeat signal from the master. If the heartbeat stops (network partition, hardware
failure, process crash), the highest-priority backup promotes itself to master, claims the VIP via a
gratuitous ARP broadcast, and begins receiving traffic. From the client's perspective, the VIP never
changes; the failover is invisible.

\texttt{Keepalived} is the de-facto Linux implementation of VRRP for this purpose. HAProxy clusters
in most on-premises deployments use a HAProxy + Keepalived combination. Failover time is typically
under one second for a pre-warmed active--passive pair --- fast enough that TCP connections in flight
may be dropped and retried, but persistent-connection applications and most retry-capable clients
survive transparently.

\begin{industrybox}{HAProxy Performance (HAProxy, Inc. 2024)}
  HAProxy running on a 64-core AWS Graviton4 instance with AWS-LC cryptographic acceleration
  achieved over 2,000,000 HTTP requests per second over TLS in benchmarks published by HAProxy,
  Inc. in 2024. On a more modest 4-vCPU instance, HAProxy sustains over 430,000 non-TLS
  connections per second with sub-millisecond latency. TLS handshake overhead is the dominant
  cost: enabling TLS 1.2 on the same hardware reduces throughput to approximately 68,000 connections
  per second due to asymmetric cryptography. These numbers illustrate why teams investing in TLS
  performance at the load balancer layer choose hardware with AES-NI acceleration or offload TLS
  to dedicated hardware security modules.
\end{industrybox}

\subsection{Cloud-Managed Load Balancers}

Cloud providers abstract away the HA problem entirely. AWS Elastic Load Balancing, Google Cloud
Load Balancing, and Azure Load Balancer are distributed services with no single master node ---
they run as a fleet of managed infrastructure with SLAs of 99.99\% or higher. You configure them
via API; you never SSH into a load balancer instance.

\begin{thinkbox}
  \textbf{Does this mean I never think about LB failure in interviews?} Not quite. Even
  cloud-managed load balancers can become a bottleneck or throttle point at extreme scale. AWS NLB
  pre-warms its capacity automatically but has per-account limits on the number of targets and
  connections. At sufficiently high scale (tens of millions of concurrent connections), engineers
  at companies like Cloudflare and Google have built custom load-balancing software
  (Cloudflare's Unimog, Google's Maglev) precisely because no off-the-shelf product handles
  their specific traffic patterns. Mentioning this in an interview signals architectural awareness.
\end{thinkbox}

% ─────────────────────────────────────────────────────────────────────────────
\section{Modern Load Balancer Landscape}

\subsection{Client-Side Load Balancing}

There is a category of load balancing that does not use a central proxy at all: \emph{client-side
load balancing}. In a microservices architecture, each service instance maintains a list of
addresses of its dependencies (fetched from a service registry like Netflix Eureka or Kubernetes
DNS). When Service A wants to call Service B, it picks a B instance itself using a local round-robin
or least-connections algorithm and connects directly. No intermediate proxy touches the request.

Netflix pioneered this pattern with the Ribbon library (now mostly superseded by Spring Cloud
LoadBalancer and gRPC's built-in client-side balancing). The advantage is zero-overhead routing
--- no extra network hop through a proxy. The disadvantage is that routing logic lives in every
service client, and health-check state must be propagated via the service registry. This pattern
is appropriate for east--west (service-to-service) traffic inside a cluster. It is \emph{not}
appropriate for north--south traffic arriving from the public internet, which still needs a
centralised ingress load balancer.

\subsection{AWS Load Balancer Family}

AWS offers three distinct managed load balancers, each targeting a different layer and use case:

\medskip
\begin{tabularx}{\linewidth}{@{}lXX@{}}
  \toprule
  \textbf{Product} & \textbf{Layer / Protocol} & \textbf{Primary Use Case} \\
  \midrule
  \textbf{ALB} (Application LB) & L7: HTTP, HTTPS, WebSocket, gRPC &
    Microservices routing, path-based rules, host-based routing, WAF integration \\
  \addlinespace
  \textbf{NLB} (Network LB) & L4: TCP, UDP, TLS &
    Ultra-low latency, non-HTTP protocols, gaming, IoT, static IP requirement \\
  \addlinespace
  \textbf{GWLB} (Gateway LB) & L3/L4 transparent pass-through &
    Inline network appliances (firewalls, IDS/IPS), packet inspection \\
  \bottomrule
\end{tabularx}
\medskip

The practical rule: reach for ALB by default for HTTP/HTTPS workloads. Switch to NLB when you need
sub-millisecond latency, static IP addresses for IP allowlisting by downstream partners, or
protocols that are not HTTP. Use GWLB only when you are inserting security appliances into a traffic
path.

% ─────────────────────────────────────────────────────────────────────────────
\section{Choosing Your Strategy: Decision Framework}

\subsection{Algorithm Selection}

\medskip
\begin{tabularx}{\linewidth}{@{}lXX@{}}
  \toprule
  \textbf{Algorithm} & \textbf{Best When} & \textbf{Avoid When} \\
  \midrule
  Round Robin & Homogeneous fleet, uniform request cost & Variable request duration \\
  \addlinespace
  Weighted RR & Mixed instance types, canary deployments & Real-time load is unpredictable \\
  \addlinespace
  Least Connections & Variable request duration, DB-heavy & Fleet is heterogeneous (biases to fast servers) \\
  \addlinespace
  Least Response Time & You want adaptive, latency-aware routing & High state-tracking overhead is a concern \\
  \addlinespace
  IP Hash & Simple affinity, no cookie support & Backend pool changes frequently \\
  \addlinespace
  Consistent Hashing & Distributed caches, database sharding, key-based routing & Simple stateless HTTP APIs \\
  \bottomrule
\end{tabularx}
\medskip

\subsection{L4 vs.\ L7 Selection}

\medskip
\begin{tabularx}{\linewidth}{@{}lXX@{}}
  \toprule
  \textbf{Need} & \textbf{Use} & \textbf{Why} \\
  \midrule
  Content-based routing (path, header) & L7 & L4 cannot read payload \\
  WebSocket or gRPC & L7 & Requires HTTP upgrade or HTTP/2 parsing \\
  Non-HTTP protocol (MySQL, Redis, SMTP) & L4 & L7 parser would not understand the protocol \\
  Ultra-low latency ($<$ 0.5ms added) & L4 & No parsing or TLS termination overhead \\
  TLS termination & L7 & L4 pass-through keeps TLS encrypted end-to-end \\
  DDoS absorption at scale & L4 first, L7 second & L4 can drop SYN floods before L7 parsing \\
  \bottomrule
\end{tabularx}
\medskip

% ─────────────────────────────────────────────────────────────────────────────
\section{Critical Numbers Reference}

\medskip
\begin{tabularx}{\linewidth}{@{}lXX@{}}
  \toprule
  \textbf{Metric} & \textbf{Value} & \textbf{Source / Context} \\
  \midrule
  HAProxy max RPS (TLS, Graviton4, 64-core) & $>$2,000,000 req/s & HAProxy, Inc. benchmark 2024 \\
  \addlinespace
  HAProxy non-TLS connections/s (4-vCPU) & $>$430,000 conn/s & HAProxy benchmark \\
  \addlinespace
  HAProxy TLS 1.2 connections/s (4-vCPU) & $\approx$68,000 conn/s & HAProxy benchmark \\
  \addlinespace
  AWS ALB health check interval (default) & 30 s & AWS documentation 2024 \\
  \addlinespace
  AWS ALB unhealthy threshold (default) & 2 failures & AWS documentation 2024 \\
  \addlinespace
  Cassandra consistent hash ring size & $2^{64}$ (Murmur3) & Cassandra 4.0+ architecture docs \\
  \addlinespace
  Cassandra virtual nodes per physical node & 256 (default) & Cassandra 4.0+ \\
  \addlinespace
  Redis Cluster hash slots & 16,384 (CRC16) & Redis Cluster specification \\
  \addlinespace
  VRRP failover time (HAProxy + Keepalived) & $<$1 second & Keepalived documentation \\
  \addlinespace
  Cloudflare PoP count & 330+ cities & Cloudflare infrastructure docs 2024 \\
  \addlinespace
  Cloudflare anycast population reach & 95\% within 50ms & Cloudflare docs 2024 \\
  \addlinespace
  Google Cloud LB availability SLA & 99.99\% & Google Cloud documentation \\
  \addlinespace
  Keys remapped when 1 node added (consistent hashing) & $1/N$ of total keys & Mathematical property \\
  \addlinespace
  Keys remapped when 1 node added (modular hashing) & $\approx$100\% of total keys & Mathematical property \\
  \bottomrule
\end{tabularx}
\medskip

% ─────────────────────────────────────────────────────────────────────────────
\section{System Design Application}

\subsection{Example 1: Designing a URL Shortener at Scale}

The interviewer asks you to design a URL shortener (think bit.ly) that handles 100,000 redirects per
second globally. The first scaling decision is where the load balancer lives and what type it is.

At the internet edge you place an L7 ALB (or an Anycast CDN-fronted solution). The LB routes on
URL path: requests to \texttt{/{shortcode}} go to the redirect service pool; requests to
\texttt{/api/shorten} go to the creation service pool. This path-based routing separates read from
write traffic, allowing you to scale each independently. Redirect traffic is read-heavy and hot, so
you scale that pool aggressively and potentially place it behind a CDN. The creation service handles
writes and can be a smaller fleet.

Health checks for the redirect service should be aggressive (5-second interval, 2 failure threshold)
because a degraded redirect service directly impacts user experience. The algorithm should be Least
Connections, since some shortcodes may be hot (a viral link) causing certain cache-miss requests
to take much longer than cache hits.

\subsection{Example 2: Designing a Real-Time Multiplayer Game Backend}

A real-time multiplayer game requires sub-millisecond game state synchronisation over UDP. Using an
L7 HTTP load balancer here would be wrong: it cannot handle UDP, and its parsing overhead would
add unacceptable latency. The correct choice is an L4 load balancer (AWS NLB in AWS environments)
configured for UDP. NLB at L4 can forward UDP packets at near-wire speed without parsing any
payload.

Within a game session, you \emph{do} want all packets from the same client to reach the same game
server (session state lives in that server's memory). IP Hash or Consistent Hashing on the client
IP achieves this without cookies. But the limitation is clear: if that game server fails, the
session is lost and the player must reconnect to a new session. For a game this is acceptable;
the game reconnects automatically. For a banking application it would not be.

\subsection{Example 3: Scaling an Existing Monolith with a Database Bottleneck}

Your interviewer says: ``The service is already behind an ALB but you're seeing CPU spikes on
individual EC2 instances and some requests timing out.'' The first question to ask is: are the
timeouts correlated with specific backends, or random? If correlated, the health-check configuration
is too lenient --- a degraded backend is staying in rotation too long. Tighten the unhealthy
threshold from 5 to 2. If the load distribution is uneven (one instance always busy, others idle),
switch from Round Robin to Least Connections to let the balancer adapt to request heterogeneity.
If the database is the bottleneck, the LB cannot fix that --- you need read replicas, a cache layer
(L1-07), or database sharding. This distinction --- \emph{what the LB can fix} versus \emph{what
it cannot} --- is exactly the kind of precision interviewers are looking for.

% ─────────────────────────────────────────────────────────────────────────────
\section{Curiosity Corner}
\addcontentsline{toc}{section}{Curiosity Corner}

\begin{curiobox}{How does Google's Maglev load balancer work, and why did Google build it?}
  Google built Maglev because commodity software load balancers could not scale to Google's traffic
  volumes while maintaining consistent hashing across a fleet of stateless forwarding nodes.
  Maglev is a distributed software load balancer: instead of one active LB and one passive standby,
  Maglev runs on hundreds of commodity Linux servers in parallel. ECMP (Equal-Cost Multi-Path)
  routing at the top-of-rack switch level distributes packets across the Maglev pool. Each Maglev
  node independently applies consistent hashing to select the same backend for a given flow ---
  without any synchronisation between nodes. This stateless, massively parallel design scales
  horizontally without limit. Google published the Maglev paper at NSDI 2016; it remains the
  canonical reference for understanding how hyperscale load balancing actually works internally.
\end{curiobox}

\begin{curiobox}{What is the difference between a load balancer and a reverse proxy?}
  The terms are often used interchangeably but describe different primary functions. A reverse proxy
  is a server that sits in front of backends and forwards client requests on their behalf, providing
  security (hiding backend IPs), SSL termination, caching, and request rewriting. A load balancer's
  primary function is distributing requests across a pool of backends to prevent any one from being
  overwhelmed. In practice, every modern L7 load balancer is also a reverse proxy (NGINX, HAProxy,
  Envoy all serve both roles simultaneously), and every reverse proxy with multiple upstream
  configurations is also a load balancer. The distinction matters mostly in vendor documentation
  and in interviews where you want to signal awareness that the two concepts are layered, not
  competing.
\end{curiobox}

\begin{curiobox}{What happens if the load balancer's session table gets corrupted or the LB restarts?}
  With sticky sessions implemented via cookies, a load balancer restart is less catastrophic than
  it might seem: the session affinity information lives \emph{in the cookie on the client}, not in
  the LB's memory. When the LB restarts, it simply reads the cookie on the next request and routes
  accordingly --- no state is lost from the LB's perspective. With IP-hash-based affinity (no
  cookie), the mapping is computed deterministically from the client IP every time, so there is
  also no state to lose on restart. The more dangerous failure mode is when the backend pool
  changes (a server is added or removed) while using IP hash: the deterministic mapping shifts and
  sessions are disrupted for clients whose hash falls on the boundary. Consistent hashing minimises
  this by design.
\end{curiobox}

\begin{curiobox}{How do load balancers handle WebSocket connections?}
  WebSocket begins as an HTTP/1.1 connection that is upgraded via the \texttt{Upgrade: websocket}
  header. An L7 load balancer must recognise this upgrade request and handle it specially: once
  the WebSocket handshake completes, the connection becomes a persistent, bidirectional TCP
  stream. The load balancer must keep this connection pinned to the same backend for its entire
  lifetime (since WebSocket state lives on that backend), and it must not impose request-level
  timeouts that would prematurely kill long-lived connections. AWS ALB supports WebSocket natively
  and handles the upgrade automatically. Discord, which maintains over 12 million concurrent
  WebSocket connections as of 2024, uses a fleet of internal L7 proxies (Elixir-based) that handle
  the upgrade and fan out events --- the public-facing layer is fronted by AWS NLB for raw TCP
  capacity, with L7 logic handled by application-level code behind it.
\end{curiobox}

\begin{curiobox}{What is ``connection draining'' and why does it matter for deployments?}
  When a backend is removed from rotation --- whether for a planned deployment, autoscale-in event,
  or health-check failure --- there may be in-flight requests that are already being processed by
  that backend. Immediately closing the connection would break those requests. \emph{Connection
  draining} (called ``deregistration delay'' in AWS terminology) is the grace period during which
  the load balancer stops sending \emph{new} requests to the departing backend while allowing
  existing in-flight requests to complete naturally, up to a configurable timeout (AWS default: 300
  seconds). Without connection draining, rolling deployments cause a percentage of requests to fail
  at each step. AWS ALB and NLB both implement connection draining. Teams with very short request
  durations (sub-100ms APIs) often reduce the deregistration delay to 15--30 seconds to speed up
  deployments without risking request failures.
\end{curiobox}

\begin{curiobox}{Can a load balancer itself become a performance bottleneck?}
  Yes, and this is a sophisticated point that distinguishes senior-level answers. A single load
  balancer instance has a maximum throughput --- in bytes per second, packets per second, and
  concurrent connections. If your application traffic exceeds that ceiling, the load balancer
  becomes the bottleneck instead of the backends. Cloud-managed load balancers (ALB, NLB) scale
  their internal capacity automatically, though AWS NLB pre-warms at predictable rates and may
  need pre-warming requests for sudden massive traffic spikes. On-premises software LBs (HAProxy,
  NGINX) are bounded by the hardware they run on. The mitigation is either to scale the LB
  horizontally (ECMP + Anycast, as Google does with Maglev) or to use DNS-level GSLB to split
  traffic across multiple independent LB clusters, each handling a geographic subset of traffic.
\end{curiobox}

% ─────────────────────────────────────────────────────────────────────────────
\newpage
\section{Self-Quiz}
\addcontentsline{toc}{section}{Self-Quiz}

\begin{quizbox}
  Answer each question as if the interviewer asked it directly. Aim for a complete answer in
  under 90 seconds per question.

  \begin{enumerate}[leftmargin=*, label=\textbf{Q\arabic*.}]

    \item \textbf{``Walk me through the difference between L4 and L7 load balancing. When would
    you use each?''}

    \item \textbf{``I've deployed my service behind a load balancer with Round Robin. Users are
    complaining that some requests are extremely slow. What might be wrong, and how would you fix
    it?''}

    \item \textbf{``Our database cluster uses consistent hashing. Explain how consistent hashing
    works and why it matters when we add a new database node.''}

    \item \textbf{``We're designing a globally distributed API. How do you route users to the
    nearest data centre, and what are the failure modes of your approach?''}

    \item \textbf{``I've heard sticky sessions are bad. But our current app stores session data
    in memory on each server. What's the right architectural fix, and how does the load balancer
    factor in?''}

  \end{enumerate}
\end{quizbox}

% ─────────────────────────────────────────────────────────────────────────────
\newpage
\section*{Model Answers}
\addcontentsline{toc}{section}{Model Answers}

\begin{answerbox}

  \textbf{A1: L4 vs.\ L7 Load Balancing.}
  An L4 load balancer routes at the TCP/UDP transport layer: it sees only IP addresses and port
  numbers, makes routing decisions from those alone, and forwards packets without reading payload.
  It is extremely fast (near wire speed), works for any protocol, and adds minimal latency. An L7
  load balancer terminates the client connection, decrypts TLS, reads the HTTP request, and routes
  based on URL path, headers, or cookies. It enables content-based routing and WebSocket support
  but adds CPU overhead and a small latency cost. I use L4 when I need ultra-low latency or a
  non-HTTP protocol (game server, raw TCP). I use L7 for HTTP microservices that need path-based
  routing, host-based virtual hosting, or gRPC.

  \medskip
  \textbf{A2: Round Robin and Slow Requests.}
  Round Robin distributes requests by count, not by cost. If some requests take 10ms and others
  take 2,000ms --- perhaps certain user IDs trigger a slow DB query --- backends that received
  several expensive requests become overloaded even though they received the same count of
  requests as their peers. The fix is to switch to Least Connections (route new requests to the
  backend with fewest active connections) or Least Response Time (track per-backend latency and
  bias toward faster backends). These algorithms are adaptive to heterogeneous request cost.

  \medskip
  \textbf{A3: Consistent Hashing.}
  Consistent hashing maps both backends and request keys (e.g., a partition key) onto a circular
  hash ring with a fixed hash space (e.g., $2^{64}$). A request key is hashed to a position on the
  ring, then routed clockwise to the nearest backend. When a new backend is added, it is placed at
  one or more positions on the ring (using virtual nodes for even distribution) and takes over only
  the ring segment between itself and its predecessor. All other keys retain their existing
  assignment. Only $1/N$ of keys are remapped on average, versus nearly 100\% with modular hashing.
  This matters enormously for a database cluster: adding a node should cause only a small fraction
  of partition keys to migrate, not a full reshuffle.

  \medskip
  \textbf{A4: Global Routing Approaches.}
  Two main approaches. GeoDNS: the authoritative DNS server returns the IP of the nearest regional
  load balancer based on the client's inferred location. Simple, but failover is limited by DNS
  TTL --- a 60-second TTL means up to 60 seconds of traffic to a dead region. Anycast: the same IP
  is announced from multiple regions via BGP; the internet routes each client to the nearest
  announcement point automatically. Failover happens in seconds as BGP reconverges, invisible to
  DNS. Failure modes of GeoDNS: stale DNS caches during regional failure. Failure modes of
  anycast: BGP routing anomalies (route leaks) can misdirect traffic globally; this is why
  anycast is paired with health-checked origin pools and automatic BGP withdrawal on failure.

  \medskip
  \textbf{A5: Fixing Server-Side Session State.}
  Sticky sessions are a band-aid, not a fix. The root problem is that session state is stored in a
  single server's memory, creating server affinity. The correct fix is to externalise session state
  to a distributed store --- Redis (for sub-millisecond reads) or DynamoDB (for durability) ---
  that any backend can read. When session state lives outside the server, the load balancer can
  route any request to any healthy backend with no affinity table. Sessions survive server failures,
  autoscaling events, and deployments without disruption. The load balancer reverts to a simple,
  stateless routing function, which is what it should be.

\end{answerbox}

% ─────────────────────────────────────────────────────────────────────────────
\newpage
\section{What's Next}
\addcontentsline{toc}{section}{What's Next}

\begin{insightbox}
  \textbf{Next: L1-05 --- Databases: Relational.}

  Load balancing distributes traffic across stateless compute. The natural next question every
  interviewer asks is: where does the state live? The answer, for most systems, begins with a
  relational database. L1-05 covers the foundations of relational storage that every system design
  answer depends on: how ACID transactions work and why they matter, how indexes are structured
  (B-trees, covering indexes, partial indexes), how replication propagates writes to read replicas,
  and how sharding partitions a monolithic database into a horizontally scalable fleet. The
  consistent hashing you learned in this lesson will reappear in L1-05 as the key distribution
  mechanism for database sharding. The latency numbers from this lesson (the cost of a round-trip
  through an LB) will anchor your intuition for how many database calls you can afford per request
  before latency SLAs are breached.

  After L1-05 and L1-06 (NoSQL), you will have covered compute routing (L1-04) and data storage
  (L1-05, L1-06) --- the two axes that define virtually every system design solution. From there,
  L1-07 (Caching) will show you how to keep the hot path out of the database entirely.
\end{insightbox}

\end{document}
